\section{Evaluation Protocol}
\label{sec:evalprotocol}

These six primitives now serve two purposes rather than one: as before, they are PLEM's
evaluation vocabulary for scoring any prediction against ground truth; as of
Section~\ref{sec:method}, three of them are also the metric each corresponds to that a
differentiable loss term is designed to approximate during training. Nothing in this section
changes to accommodate that --- the metrics stay exactly as originally defined and validated;
Section~\ref{sec:method} does the work of relating a \emph{separate} differentiable formulation
back to each one.

\textbf{Shared notation.} Let $P$ and $G$ be $H\times W$ integer label maps (prediction and
ground truth) over the class set $C = \{0, 1, 2, [3]\}$. For class $c$, define binary masks
$P_c = (P = c)$ and $G_c = (G = c)$. Every primitive in this section is applied per-class and then
reduced (macro or GT-pixel-count-weighted mean) across the classes relevant to it.

\textbf{Shared edge-case convention.} Every primitive below implements the same rule, verified
identical across \texttt{dtaf1.py:43-51}, \texttt{cldice.py:51-56}, \texttt{boundary\_f1.py:68-72},
\texttt{point\_f1.py:109-120}, and \texttt{apls.py:141-144}:
both $P_c$ and $G_c$ empty $\Rightarrow$ score $= 1.0$ (vacuously correct); exactly one of $P_c,
G_c$ empty $\Rightarrow$ score $= 0.0$ (total failure). Stated once here as a cross-cutting design
decision rather than repeated per metric.

\subsection{DTAF1 --- Distance-Tolerant, per-class F1}
\label{subsec:dtaf1}

Tolerance-radius matching via Euclidean distance transforms
(\texttt{scipy.ndimage.distance\_transform\_edt}):
\begin{align}
  \mathrm{TP}_{\mathrm{pred}} &= |\{x \in P_c : \mathrm{EDT}(1-G_c)(x) \le d_c\}| \\
  \mathrm{TP}_{\mathrm{gt}}   &= |\{x \in G_c : \mathrm{EDT}(1-P_c)(x) \le d_c\}| \\
  \mathrm{Precision}_c &= \mathrm{TP}_{\mathrm{pred}} / |P_c|, \quad
  \mathrm{Recall}_c = \mathrm{TP}_{\mathrm{gt}} / |G_c| \\
  F1_c &= \frac{2 \cdot \mathrm{Precision}_c \cdot \mathrm{Recall}_c}
               {\mathrm{Precision}_c + \mathrm{Recall}_c}
\end{align}
Reduction across classes (both are always computed; \texttt{reduction} just selects which is
surfaced as the top-level \texttt{dtaf1} key):
\begin{align}
  \mathrm{DTAF1}_{\mathrm{macro}} &= \tfrac{1}{|C|}\textstyle\sum_c F1_c \\
  \mathrm{DTAF1}_{\mathrm{weighted}} &= \left(\textstyle\sum_c n_c F1_c\right) \big/
    \left(\textstyle\sum_c n_c\right), \quad n_c = |G_c|
\end{align}

\begin{lstlisting}[language=Python, caption={Tolerance-matching core, \texttt{metrics/dtaf1.py:53-64}}]
dist_from_gt = distance_transform_edt(1 - g)
dist_from_pred = distance_transform_edt(1 - p)
tp_pred = int(((p == 1) & (dist_from_gt <= d)).sum())
tp_gt = int(((g == 1) & (dist_from_pred <= d)).sum())
precision = tp_pred / n_pred
recall = tp_gt / n_gt
f1 = (2 * precision * recall / (precision + recall)
      if (precision + recall) > 0 else 0.0)
\end{lstlisting}

\texttt{DEFAULT\_TOLERANCES = \{"road": 10, "building": 2\}} (pixels; see
Section~\ref{subsec:tolerance-radii} for how these were chosen). Function signature:
\texttt{dtaf1(pred, gt, class\_config, reduction="macro") -> dict}, returning \texttt{dtaf1},
\texttt{dtaf1\_macro}, \texttt{dtaf1\_weighted}, \texttt{per\_class}. \textbf{Class-agnostic by
construction} --- scoring a new class costs zero code changes, just an added
\texttt{class\_config} entry. Section~\ref{subsec:tolerance-loss}'s \texttt{ToleranceBandLoss}
reuses this exact \texttt{class\_config} shape.

\subsection{clDice --- Centerline Dice}
\label{subsec:cldice}

(Cite Shit et~al., CVPR 2021 \cite{shit2021cldice} in Section~\ref{sec:related} --- not here.)
\begin{align}
  S_p &= \mathrm{skeletonize}(P_c), \quad S_g = \mathrm{skeletonize}(G_c) \\
  T_{\mathrm{prec}} &= |S_p \cap G_c| / |S_p|, \quad T_{\mathrm{sens}} = |S_g \cap P_c| / |S_g| \\
  \mathrm{clDice} &= \frac{2 \cdot T_{\mathrm{prec}} \cdot T_{\mathrm{sens}}}
                          {T_{\mathrm{prec}} + T_{\mathrm{sens}}}
\end{align}
Width-insensitive (a predicted road $3.3\times$ too thick still scores $\mathrm{clDice}=1.000$,
Table~\ref{tab:failure-modes} row 2) but offset-brittle: skeletons of the exact same road shifted
5px share \emph{zero} overlap ($\mathrm{clDice}=0.000$, row 1) --- this offset-sensitivity is
exactly what motivates DTAF1's distance-tolerant matching instead of a raw skeleton intersection.
Section~\ref{subsec:cldice-loss}'s \texttt{SoftClDiceLoss} reimplements this width-insensitivity
property differentiably, with a caveat (Sections~\ref{subsec:cldice-loss},~\ref{sec:discussion}).

\subsection{Boundary F1 (BF) and the IoU Baseline}
\label{subsec:bf}

(Cite Csurka et~al., BMVC 2013 \cite{csurka2013evaluation} in Section~\ref{sec:related} --- not
here.)
\begin{align}
  \mathrm{Precision} &= |\{x \in B_p : \mathrm{dist}(x, B_g) \le \tau\}| / |B_p| \\
  \mathrm{Recall} &= |\{x \in B_g : \mathrm{dist}(x, B_p) \le \tau\}| / |B_g| \\
  \mathrm{BF} &= \frac{2 \cdot \mathrm{Precision} \cdot \mathrm{Recall}}
                      {\mathrm{Precision} + \mathrm{Recall}}
\end{align}
where $B_p = \mathrm{boundary}(P_c)$, $B_g = \mathrm{boundary}(G_c)$ via morphological boundary
extraction. \texttt{iou()} is introduced here explicitly as the pixel-overlap baseline comparator
used throughout Sections~\ref{sec:motivation} and~\ref{sec:expsetup}. Section~\ref{subsec:boundary-loss}'s
\texttt{SoftBoundaryLoss} reuses Section~\ref{subsec:tolerance-loss}'s precision/recall primitive
applied to a soft boundary map instead of the full class mask, making it the most directly-shared
piece of machinery between any two \texttt{losses/} terms.

\subsection{Point F1 --- Hungarian Instance Matching}
\label{subsec:pointf1}

For the optional point class (trees, lamp posts, manhole covers): connected-component blobs are
reduced to centroids (\texttt{scipy.ndimage.label} + \texttt{center\_of\_mass}), then matched via
tolerance-restricted Hungarian assignment (\texttt{scipy.optimize.linear\_sum\_assignment}), not
greedy nearest-neighbor.

\begin{lstlisting}[language=Python, caption={Hungarian assignment core, \texttt{metrics/point\_f1.py:58-65}}]
dist = cdist(gt_c, pred_c)
big = tolerance * 1e6 + 1.0
cost = np.where(dist <= tolerance, dist, big)
row_ind, col_ind = linear_sum_assignment(cost)
valid = dist[row_ind, col_ind] <= tolerance
return int(valid.sum())
\end{lstlisting}

Justified by the clustered-points adversarial case in \texttt{TestPointInstanceMatching}: greedy
nearest-neighbor gets $\mathrm{TP}=1$ where Hungarian finds the globally optimal $\mathrm{TP}=2$
(formalized in Appendix~\ref{app:proofs}, Section~\ref{app:a6}). \textbf{This is the one primitive
with no differentiable analog in Section~\ref{sec:method}} --- Hungarian assignment and
connected-component labeling have no useful gradient path back to per-pixel logits;
Section~\ref{subsec:heatmap-loss} reformulates point supervision entirely as heatmap regression
instead, a genuine train/eval formulation mismatch discussed explicitly in
Section~\ref{sec:discussion}.

\subsection{APLS --- Raster-Skeleton Average Path Length Similarity}
\label{subsec:apls}

No vector road graph survives anywhere in this pipeline --- \texttt{datasets/common.py::rasterize\_lines}
collapses roads straight to a flat pixel mask, discarding segment/graph structure --- so APLS is
approximated entirely from raster skeletons (\texttt{metrics/apls.py}): skeletonize pred/GT masks,
build an 8-connected pixel-adjacency graph (\texttt{networkx}; edge weight = Euclidean pixel
distance), then for sampled GT control-point pairs $(u,v)$ within one GT connected component,
compare GT shortest-path length $L_{\mathrm{gt}}(u,v)$ against the predicted shortest-path length
between the nearest snap-matched predicted-graph nodes (\texttt{scipy.spatial.cKDTree}), scoring
$\max(0, 1 - |L_{\mathrm{pred}} - L_{\mathrm{gt}}| / L_{\mathrm{gt}})$ if matched, else 0.
Control-point pairs are sampled \textbf{grouped by source}, so one
\texttt{nx.single\_source\_dijkstra\_path\_length} call amortizes across many targets instead of
one independent Dijkstra call per pair --- see complexity analysis in
Appendix~\ref{app:complexity}, Section~\ref{app:b6}. This is deliberately sensitive to a failure
DTAF1's per-pixel tolerance matching misses entirely: once the skeleton fragments, shortest paths
between control points either lengthen sharply or vanish, so APLS collapses where DTAF1 does not.
\textbf{Graph shortest-path recomputation has no tractable lightweight differentiable relaxation}
(Section~\ref{sec:method}'s design assessment, Section~\ref{sec:discussion}) --- APLS and
DTAF1-Topo (Section~\ref{sec:topological}) stay eval-only permanently, a stated design decision
rather than a to-do.

\subsection{CBHM --- Composite clDice/Boundary-F1 Harmonic Mean}
\label{subsec:cbhm}

\begin{align}
  \mathrm{clDice}_{\mathrm{mean}} &= \mathrm{mean}(\mathrm{clDice}_c : c \in \mathrm{linear\_classes}) \\
  \mathrm{BF}_{\mathrm{mean}} &= \mathrm{mean}(\mathrm{BF}_c : c \in \mathrm{polygon\_classes}) \\
  \mathrm{CBHM} &= \frac{2 \cdot \mathrm{clDice}_{\mathrm{mean}} \cdot \mathrm{BF}_{\mathrm{mean}}}
                        {\mathrm{clDice}_{\mathrm{mean}} + \mathrm{BF}_{\mathrm{mean}}}
\end{align}
The harmonic mean is a \textbf{deliberate design choice}: a single class scoring 0 collapses the
whole composite to 0. Positioned explicitly as the ``harsh'' foil to DTAF1's more lenient macro
average (Table~\ref{tab:failure-modes}'s ``Road shifted 5px'' row: $\mathrm{DTAF1}=1.000$ but
$\mathrm{CBHM}=0.000$, on the \emph{same} input).

\textbf{Composition note.} These six primitives compose into the two ``final iteration''
composites covered later rather than here --- \texttt{cbhm\_soft} in Section~\ref{sec:composite}
and DTAF1-Topo in Section~\ref{sec:topological} --- and, as of this revision, into the three
geometry-aware \texttt{losses/} terms in Section~\ref{sec:method}.
